
\chapter{Screen scraping (2 of 2)}
\label{ch:screenScraping2}

\index{the easy case}
\index{screen scraping}
\index{worst of all possible worlds}

\section{The harder case: picking through the HTML tree}

Chapter~\ref{ch:screenScraping1} was the ``easy'' screen scraping case: if the
data on the web page that you want to import is in \texttt{<table>} rows, you
can (relatively) simply use Pandas' \texttt{read\_html()} function to suck it
in to Python.

For anything else, you'll have to comb through the HTML structure, figuring out
how to surgically extract what you seek. It's a major pain, but there is a
Python library by the outrageous name ``BeautifulSoup'' which makes it just a
little bit less painful. This chapter contains an overview of how to work with
it.


\section{Web Developer Tools}

\index{web developer tools}
\index{element (HTML)}

First, let me say that Firefox, Chrome, and several other browsers come
equipped with very sophisticated ``web developer tools'' that let you browse
the HTML, CSS, and rendered page, and figure out what elements got rendered via
which portions of the screen. They're very handy, and when screen scraping you
should get into the habit of using them. In Firefox, they're available via the
``Tools | Web Developer | Toggle tools...'' menu option. In Chrome, you go to
the little three vertical dots, then choose ``More Tools | Developer Tools.''


\section{Parsing with BeautifulSoup}

\index{BeautifulSoup}
\index{web crawler}
\index{URL}

As I mentioned on p.~\pageref{downloading}, it's generally recommended to
\textit{download} the web page you want to parse rather than having Python read
it for you over the network. This eliminates one more possible source of error.
(Exception: if you need a Python program to act as a \textbf{web crawler} and
read pages on demand as it discovers their URLs. More on that later.) Note that
you can still open the page in your browser even after downloading it: just go
to ``File | Open File...'' or the equivalent operation in your browser.

Once you've done this, and now have your HTML in a file called (say)
``\texttt{myfile.htm}'', you can then parse it with BeautifulSoup using:

\index{bs4@\texttt{bs4} (BeautifulSoup library)}

\begin{Verbatim}[fontsize=\small,samepage=true,frame=single,framesep=3mm]
from bs4 import BeautifulSoup
soup = BeautifulSoup(open("myfile.htm"), "html.parser")
\end{Verbatim}

\index{soup@\texttt{soup}}
\index{html.parser@\texttt{html.parser} (BeautifulSoup)}

By convention, the variable for this parsed structure is called
``\texttt{soup}''. The second argument to the \texttt{BeautifulSoup} function
tells it which \textbf{parser} we want to use to analyze the code; there are
other choices here, but stick with \texttt{"html.parser"} for now.

One possibly useful command (though limited if your web page is large) is:

\index{pretty printing@``pretty printing''}
\index{prettify@\texttt{.prettify()} (BeautifulSoup)}

\begin{Verbatim}[fontsize=\small,samepage=true,frame=single,framesep=3mm]
print(soup.prettify())
\end{Verbatim}

This ``pretty prints'' out the entire contents of the page, but in a way that
is properly indented according to the HTML's structure, rather than how it was
indented in the raw page. You may have noticed that the HTML file may or may
not be physically indented and spaced in a way reflective of the structure. The
BeautifulSoup parser is smart enough to discern the implied structure from the
tag sequence, and thus prints it how it will ``really'' be seen by a browser.

\index{top-down navigation}
\index{bottom-up navigation}

At this point, I generally take one of two approaches: either I navigate the
structure \textbf{top-down} looking for where the data I'm interested in is
situated, or I find the data by searching through the file, and use a
\textbf{bottom-up} approach to get it from there. You typically have to combine
the two approaches to get your screen-scraping dirty work done.

\subsection{The top-down approach}

\index{root node (of an HTML file)}
\index{contents@\texttt{.contents} (BeautifulSoup)}

By using ``\texttt{soup.contents}'', you get the immediate descendants of
the root node of the HTML hierarchy. Each of these could be examined for its
type, like this:

\index{type@\texttt{type()}}

\begin{Verbatim}[fontsize=\small,samepage=true,frame=single,framesep=3mm]
for child in soup.contents:
    print(type(child))
\end{Verbatim}
\vspace{-.2in}

\begin{Verbatim}[fontsize=\small,samepage=true,frame=leftline,framesep=5mm,framerule=1mm]
bs4.element.Tag
bs4.element.NavigableString
\end{Verbatim}

\index{NavigableString@\texttt{NavigableString} (BeautifulSoup)}
\index{tag (HTML)}

So our entire document is composed of an element (misleadingly named ``Tag''
here, but w/e) and a ``\texttt{NavigableString}'' which essentially just means
``string of literal text''. For elements, it's useful to print out their
``\texttt{.name}''. For strings, it's useful to print them out verbatim. Here,
we see:

\begin{Verbatim}[fontsize=\small,samepage=true,frame=single,framesep=3mm]
print(soup.contents[0].name)
print(soup.contents[1])
\end{Verbatim}
\vspace{-.2in}

\begin{Verbatim}[fontsize=\small,samepage=true,frame=leftline,framesep=5mm,framerule=1mm]
'html'
'\n'
\end{Verbatim}

\index{html@\texttt{<html>} (HTML tag)}

So the first child is the entire \texttt{html} element in all its embedded
glory, and this is followed by a single, pesky newline. Now you know.

Drilling down another level, we see:

\begin{Verbatim}[fontsize=\small,samepage=true,frame=single,framesep=3mm]
html_children = soup.contents[0]
for child in html_children:
    print(type(child))
\end{Verbatim}
\vspace{-.2in}

\begin{Verbatim}[fontsize=\small,samepage=true,frame=leftline,framesep=5mm,framerule=1mm]
bs4.element.NavigableString
bs4.element.Tag
bs4.element.NavigableString
bs4.element.Tag
bs4.element.NavigableString
\end{Verbatim}

Two tags, with three strings interleaved. Let's examine them:

\begin{Verbatim}[fontsize=\small,samepage=true,frame=single,framesep=3mm]
print(html_children.contents[0])
print(html_children.contents[1].name)
print(html_children.contents[2])
print(html_children.contents[3].name)
print(html_children.contents[4])
\end{Verbatim}
\vspace{-.2in}

\begin{Verbatim}[fontsize=\small,samepage=true,frame=leftline,framesep=5mm,framerule=1mm]
'\n'
'head'
'\n'
'body'
'\n'
\end{Verbatim}

You're starting to see a pattern: a lot of the strings are just newline
characters that split up the HTML text file and make it more readable. What we
really care about is the content under the tags. Let's dig one deeper:

\begin{Verbatim}[fontsize=\small,samepage=true,frame=single,framesep=3mm]
head_children = html_children.contents[1].contents
print(head_children[0])
print(head_children[1].name)
print(head_children[2])
print(head_children[3].name)
print(head_children[4])
\end{Verbatim}
\vspace{-.2in}

\begin{Verbatim}[fontsize=\small,samepage=true,frame=leftline,framesep=5mm,framerule=1mm]
'\n'
'title'
'\n'
'style'
'\n'
\end{Verbatim}

Okay, you get the idea. We can drill down to whatever level we want just by
following the appropriate children. To get the title of the page, we could do:

\index{title@\texttt{<title>} (HTML tag)}

\begin{Verbatim}[fontsize=\small,samepage=true,frame=single,framesep=3mm]
print(head_children[1].contents[0])
\end{Verbatim}
\vspace{-.2in}

\begin{Verbatim}[fontsize=\small,samepage=true,frame=leftline,framesep=5mm,framerule=1mm]
'All About Stephen'
\end{Verbatim}

\index{title@\texttt{<title>} (HTML tag)}
\index{head@\texttt{<head>} (HTML tag)}

(This is because the second element (element \#1) of the \texttt{<head>} element
is the \texttt{<title>} element, and its only child is the string
\texttt{\textquotesingle All About Stephen\textquotesingle}.)

Sometimes a quicker method to find your way through the forest is to use this
syntax\footnote{That thing in the boxies -- ``\texttt{[ x.name for x in
body\_children ]}'' -- is called a ``\textbf{list comprehension},'' by the way,
and is an advanced Python technique that is sometimes very useful.} at the
interactive console:

\index{list!comprehension}

\begin{Verbatim}[fontsize=\small,samepage=true,frame=single,framesep=3mm]
print([ x.name for x in body_children ])
\end{Verbatim}
\vspace{-.2in}

\begin{Verbatim}[fontsize=\small,samepage=true,frame=leftline,framesep=5mm,framerule=1mm]
[None, 'h1', None, 'p', None, 'p', None, 'h2', None, 'table', None]
\end{Verbatim}

\index{p@\texttt{<p>} (HTML tag)}

You can better see at a glance what's what. Now that we know the first
``\texttt{p}'' (paragraph) is element \#3 of the body, we could find that
hyperlink via:

\begin{Verbatim}[fontsize=\small,samepage=true,frame=single,framesep=3mm]
print([ x.name for x in list(body_children)[3] ])
\end{Verbatim}
\vspace{-.2in}

\begin{Verbatim}[fontsize=\small,samepage=true,frame=leftline,framesep=5mm,framerule=1mm]
[ None, 'a', None, 'b', None, 'u', None]
\end{Verbatim}

and grab it via:

\begin{Verbatim}[fontsize=\small,samepage=true,frame=single,framesep=3mm]
dat_hyperlink =  list(list(body_children)[3])[1]
\end{Verbatim}

\index{dictionary}
\index{attribute!of an HTML element}
\label{dictSyntax}

BeautifulSoup lets you access the attributes of an HTML element by treating
the element as a dictionary. Pretty cool. We can snag the actual link, then,
with:

\begin{Verbatim}[fontsize=\small,samepage=true,frame=single,framesep=3mm]
print(dat_hyperlink['href'])
\end{Verbatim}
\vspace{-.2in}

\begin{Verbatim}[fontsize=\small,samepage=true,frame=leftline,framesep=5mm,framerule=1mm]
'http://buildings.umwblogs.org/james-farmer-hall/'
\end{Verbatim}

\subsection{The bottom-up approach}

Sometimes, when you try to dive from the top of the document all the way to
what you want, you can find yourself in a tangled mess. In these cases it's
easier to search for the data you want using BeautifulSoup's search facilities.

Let's illustrate this by creating a simple \texttt{DataFrame} that contains
each role and its percentage allocation of Stephen's time. We want \texttt{df}
to look something like this:

\begin{Verbatim}[fontsize=\small,samepage=true,frame=leftline,framesep=5mm,framerule=1mm]
                role  perc
0  College professor  50.0
1                Dad  30.0
           . . .
\end{Verbatim}

To find our way through the forest, we search through the HTML with developer
tools, locating one of the pieces of data we know is there -- for instance, the
number \texttt{50}, which we can see on the page is one of the percentages.
When we click on this (or just search for it in a text file), we find ourselves
in the middle of the HTML content in Figure~\ref{fig:bottomUpSearch}.

\begin{figure}[ht]
\centering
\begin{Verbatim}[fontsize=\footnotesize,samepage=true,frame=lines]
...
<table border=2 cellpadding=15>
<tr><td class="pic"><img src="a.jpg"></td>
    <td class="role">College professor</td>
    <td class="description">
    Stephen teaches computer science and data science to willing undergraduates,
    shaping them into the heroes of tomorrow.
    </td>
    <td>
    <ul>
    <li>Super rewarding</li>
    <li>Exciting</li>
    <li>Ever-changing</li>
    <li>Fun</li>
    </ul>
    </td>
    <td>
    <b>50%</b> of total time
    </td>
</tr>
...
\end{Verbatim}
\caption{Using a bottom-up approach -- locating the text ``\texttt{50}'' in the
HTML file (stare towards the bottom of that HTML text) and getting one's
bearings regarding where it is in relation to other elements.}
\label{fig:bottomUpSearch}
\end{figure}
\normalsize

\index{td@\texttt{<td>} (HTML tag)}
\index{b@\texttt{<b>} (HTML tag)}
\index{class@\texttt{class} (HTML/CSS attribute)}

Looks like the role name we want is the text under a \texttt{td} element whose
class\footnote{In HTML, a ``\textbf{class}'' is an attribute used to classify
certain elements to be picked up by CSS formatting instructions. If this were a
Web development textbook, it would talk about when and how to use
``\texttt{class}''. But for our purposes, it's just one of many handy things we
might hook on to in order to find our data in the morass.} is
``\texttt{role}''. And it looks like the percentage we want is (*groan*) in a
boldface (\texttt{b}) tag, underneath the last \texttt{td} tag of the row,
which has no \texttt{class}.

\index{find\_all@\texttt{.find\_all()} (BeautifulSoup)}

BeautifulSoup's \texttt{.find\_all()} method makes this slightly less of a
headache. To quickly grab all the \texttt{tr} elements in the document, we can
do this:

\begin{Verbatim}[fontsize=\small,samepage=true,frame=single,framesep=3mm]
trs = list(soup.find_all("tr"))
\end{Verbatim}

\index{td@\texttt{<td>} (HTML tag)}

The variable \texttt{trs} is now a list of \textit{elements} (BeautifulSoup
``\texttt{Tag}'' objects) each of which we can examine for (say) \texttt{td}
children with a \texttt{role} \texttt{class} attribute. And you can call
\texttt{.find\_all()} on \textit{any} tag (to search from that point in the
hierarchy down), not just on the top-level \texttt{soup} element. So:

\index{in@\texttt{in}}
\index{array@\texttt{array()} (NumPy)}
\index{dtype@\texttt{dtype}}

\begin{Verbatim}[fontsize=\small,samepage=true,frame=single,framesep=3mm]
roles = np.array([], dtype=object)
trs = soup.find_all("tr")
for tr in trs:
    tds = tr.find_all("td")
    for td in tds:
        if "class" in td.attrs and "role" in td["class"]:
            ...
\end{Verbatim}

\index{list}
\index{dictionary}
\index{key-value pair}
\index{attribute!of an HTML element}
\index{attrs@\texttt{.attrs} (BeautifulSoup)}

The ``\texttt{.attrs}'' syntax gives us a \textit{dictionary} of key-value
pairs which represent the attributes of an element. So by saying
``\texttt{"class" in td.attrs}'' we're first checking whether there is a
\texttt{class} attribute in the current \texttt{td} element at \textit{all}
before then attempting to get its value using the dictionary syntax mentioned
on p.~\pageref{dictSyntax}.

\label{classIsAttribute}

Btw, one slight surprise in the above \texttt{if} statement might be that we're
checking whether \texttt{"role"} is \textit{\texttt{in}} \texttt{td["class"]}
instead of simply being \texttt{==} to \texttt{td["class"]}. That's because in
HTML, the \texttt{class} attribute of an tag can have \textit{multiple}
(space-separated) values, not just one, and so \texttt{td["class"]} is actually
a \textit{list}, not a single string.

\index{get\_text@\texttt{.get\_text()} (BeautifulSoup)}

Okay, now when we reach the body of that \texttt{if} statement, we have
successfully found one of the desired \texttt{td} attributes. How do we then
get the text under it? Just call \texttt{.get\_text()} on it, and then append
it to the list. This creates our desired list of roles that we can stick in our
\texttt{DataFrame}.

\begin{Verbatim}[fontsize=\small,samepage=true,frame=single,framesep=3mm]
            ...
            roles.append(td.get_text())
\end{Verbatim}

To get the percentages, we just need a slightly different approach: first,
we're going to assume the percentage is in the \textit{last} \texttt{td} tag,
and second, we need to go inside that to retrieve the percentage from the
embedded boldface tag. Adding this to our \texttt{for} loop almost gets us
there:

\begin{Verbatim}[fontsize=\small,samepage=true,frame=single,framesep=3mm]
percs = []
...
    the_bold = tds[-1].find("b")
    percs.append(the_bold.get_text())
\end{Verbatim}

although we'd like to also strip that percentage sign off and convert it to an
actual numeric value, so we'll do this instead:

\begin{Verbatim}[fontsize=\small,samepage=true,frame=single,framesep=3mm]
percs = []
...
    the_bold = tds[-1].find("b")
    percs.append(float(the_bold.get_text()[:-1]))
\end{Verbatim}

\index{find@\texttt{.find()} (BeautifulSoup)}
\index{find\_all@\texttt{.find\_all()} (BeautifulSoup)}

and now we can create our DataFrame with both columns. (Note that
``\texttt{.find()}'' does the same thing as ``\texttt{.find\_all()}'' except
that it only grabs the first matching element rather than all of them.) The
complete code is shown in Figure~\ref{fig:getRolesComplete}.

\begin{figure}[ht]
\begin{Verbatim}[fontsize=\small,samepage=true,frame=single,framesep=3mm]
roles = []
percs = []

trs = soup.find_all("tr")
for tr in trs:
    tds = tr.find_all("td")
    for td in tds:
        if 'class' in td.attrs and 'role' in td['class']:
            roles.append(td.get_text())
    the_bold = tds[-1].find("b")
    percs.append(float(the_bold.get_text()[:-1]))

df = pd.DataFrame({'role':roles, 'perc':percs})[['role','perc']]
\end{Verbatim}
\caption{The complete code example to screen scrape the roles and percentages
from the web page in Figure~\ref{fig:allAboutStephen}
(p.~\pageref{fig:allAboutStephen}).}
\label{fig:getRolesComplete}
\end{figure}

\subsection{Other BeautifulSoup features}

The BeautifulSoup library has a myriad of alternatives for navigating and
searching the tree. Here are a few.

\subsubsection{Searching by id or class}

\index{find@\texttt{.find()} (BeautifulSoup)}
\index{find\_all@\texttt{.find\_all()} (BeautifulSoup)}
\index{id@\texttt{id} (BeautifulSoup)}
\index{class\_@\texttt{class\_} (BeautifulSoup)}

\texttt{.find()} and \texttt{.find\_all()} each have optional \texttt{id} and
\texttt{class\_} (note the trailing underscore!) parameters that can be used
to restrict a search to only elements with matching attributes.

\begin{Verbatim}[fontsize=\small,samepage=true,frame=single,framesep=3mm]
soup.find_all("td", class_="pic")
\end{Verbatim}

\subsubsection{Using the Python ``dot'' syntax}

If we know the names of the tags, we can navigate a structure by using a
Pandas-like ``dot'' syntax. Check it out:

\begin{Verbatim}[fontsize=\small,samepage=true,frame=single,framesep=3mm]
print(soup.head.title)
print(soup.head.name)
print(soup.head.title.string)
print(soup.body.table.td.img['src'])
\end{Verbatim}
\vspace{-.2in}

\begin{Verbatim}[fontsize=\small,samepage=true,frame=leftline,framesep=5mm,framerule=1mm]
<title>All About Stephen</title>
title
'All About Stephen'
'a.jpg'
\end{Verbatim}


\subsubsection{Accessing all attributes}

\index{tag (HTML)}
\index{attrs@\texttt{.attrs} (BeautifulSoup)}
\index{dictionary}

As we've seen, in addition to treating an element (``\texttt{Tag}'') as a
dictionary to retrieve its element values, you can also use the
``\texttt{.attrs}'' syntax to get the dictionary directly:

\begin{Verbatim}[fontsize=\small,samepage=true,frame=single,framesep=3mm]
print(soup.body.find("a").attrs)
print(soup.body.find("td").attrs)
\end{Verbatim}
\vspace{-.2in}

\begin{Verbatim}[fontsize=\small,samepage=true,frame=leftline,framesep=5mm,framerule=1mm]
{'href': 'http://buildings.umwblogs.org/james-farmer-hall/'}
{'class': ['pic']}
\end{Verbatim}

(As noted on p.~\pageref{classIsAttribute}, see how the value of the
``\texttt{class}'' attribute is a list here.)

\subsubsection{Getting just text with \texttt{.get\_text()}}

\index{get\_text@\texttt{.get\_text()} (BeautifulSoup)}

The aforementioned \texttt{.get\_text()} method is actually a bit more
powerful than just extracting the contents of a single plain text element
underneath an element. It can also be used to get \textit{all the text} from
that point down in the hierarchy, ignoring all tags. Check it out:

\begin{Verbatim}[fontsize=\small,samepage=true,frame=single,framesep=3mm]
soup.body.find("p").get_text()
\end{Verbatim}
\vspace{-.2in}

\begin{Verbatim}[fontsize=\small,samepage=true,frame=leftline,framesep=5mm,framerule=1mm]
'Deep in the bowels of Farmer Hall resides a\ncarbon-based life form known to
outsiders only as Stephen. This\nmysterious being fulfills several roles not
widely known to\nhumankind.'
\end{Verbatim}

\subsubsection{Finding more than one tag name at once}

\index{find\_all@\texttt{.find\_all()} (BeautifulSoup)}
\index{list}

You can pass a \textit{list} to \texttt{.find\_all()} to match more than one
type of tag in a search:

\begin{Verbatim}[fontsize=\small,samepage=true,frame=single,framesep=3mm]
soup.find_all(["td", "th", "p"])
\end{Verbatim}

\index{regex (regular expression)}

And if you know how to use \textbf{regular expressions} (outside the scope of
this class) you can use them as arguments to \texttt{.find\_all()} and
friends:

\begin{Verbatim}[fontsize=\small,samepage=true,frame=single,framesep=3mm]
import re
soup.find_all(re.compile("^q"))
soup.find_all(re.compile("r"))
soup.find_all(re.compile("s$"))
\end{Verbatim}

(finds all tags whose name begin with \texttt{q}, contain \texttt{r}, or end
with \texttt{s}, respectively.)

\subsubsection{CSS selectors}

\index{CSS}
\index{select@\texttt{.select()} (BeautifulSoup)}

Finally, if you're familiar with CSS selector syntax, you can do all that jazz
by calling the \texttt{.select()} method:

\begin{Verbatim}[fontsize=\small,samepage=true,frame=single,framesep=3mm]
soup.select("div#master, tr td.data")
\end{Verbatim}

